% Work/MAT — 1. ročník, hodina 61.
% Návaznost: 59 Co je funkce; 60 Definiční obor a obor hodnot.
% Zdroje: ucebnice/02_Funkce_a_rovnice/Zaklady_funkci/Funkce.tex,
% části Monotónnost, Sudá a lichá funkce, Prostá funkce.
% Metodicky: Realisticky.cz, zejména 020113 (rostoucí a klesající funkce).
% Příklady a grafy jsou vlastní. Průsečíky necháváme konkrétním funkcím.
\documentclass[11pt,a4paper]{article}
% Spolecny styl pro ucitelske pripravy MAT (Work/MAT)
% Verze 1.1, 2026-10-23
\NeedsTeXFormat{LaTeX2e}
\ProvidesPackage{priprava}[2026/10/23 v1.1 Sablona ucitelskych priprav MAT]

% Zakladni typografie a cestina
\RequirePackage[T1]{fontenc}
\RequirePackage[utf8]{inputenc}
\RequirePackage[czech]{babel}
\RequirePackage{lmodern}
\RequirePackage{microtype}
\RequirePackage{amsmath,amssymb}
\RequirePackage{xcolor}
\RequirePackage{enumitem}
\RequirePackage[most]{tcolorbox}
\RequirePackage{tabularx,booktabs,array}
\RequirePackage{fancyhdr}
\RequirePackage{needspace}
\RequirePackage{geometry}
\RequirePackage[hidelinks]{hyperref}

\geometry{a4paper,top=17mm,bottom=18mm,left=23mm,right=23mm}
\definecolor{notegray}{gray}{0.39}
\setlength{\parindent}{0pt}
\setlength{\parskip}{2.5pt plus 1pt minus 0.5pt}
\setlist[itemize]{topsep=2pt,itemsep=1pt,parsep=0pt,leftmargin=1.45em}
\setlist[enumerate]{topsep=2pt,itemsep=2pt,parsep=0pt,leftmargin=1.65em}
\setcounter{secnumdepth}{0}
\renewcommand{\arraystretch}{1.12}

% Zahlavi: vlevo rocnik | globalni tema, vpravo cislo: nazev hodiny
% Zapati zustava zcela prazdne.
\pagestyle{fancy}
\fancyhf{}
\renewcommand{\headrulewidth}{0.35pt}
\renewcommand{\footrulewidth}{0pt}
\setlength{\headheight}{14pt}
\newcommand{\nastavPripravu}[4]{%
  \fancyhead[L]{\small #1\enspace\textbar\enspace #2}%
  \fancyhead[R]{\small #3: #4}%
}

% Minimalisticke cernobile boxy. Kazda definice zvlast.
\tcbset{commonlesson/.style={enhanced,breakable=false,colback=white,colframe=black,
  colbacktitle=white,coltitle=black,boxrule=0.45pt,
  arc=0.8mm,outer arc=0.8mm,left=3mm,right=3mm,top=1.6mm,bottom=1.6mm,
  toptitle=1mm,bottomtitle=0.7mm,fonttitle=\bfseries,
  attach boxed title to top left={xshift=2.5mm,yshift=-1.15mm},
  boxed title style={colback=white,colframe=white,boxrule=0pt,left=1mm,right=1mm}}}
\newtcolorbox{zadaniBox}[1][]{commonlesson,title={Zadání},#1}
\newtcolorbox{prikladBox}[1][]{commonlesson,title={Příklad},#1}
\newtcolorbox{definiceBox}[1][]{commonlesson,title={Definice},#1}

% Spolecny box pro zadani a jeho reseni, bez mezery mezi castmi.
% Pouziti:
%   \begin{prikladReseniBox}
%     Zadani ulohy.
%     \tcblower
%     \begin{align*} ... \end{align*}
%   \end{prikladReseniBox}
% Spodni cast ma jemny sedy podklad a delici linku; lze v ni sazet
% vice radku matematiky. Obe casti maji jen jeden spolecny ramecek.
\tcbset{solutionlesson/.style={
  commonlesson,
  enhanced jigsaw,
  breakable,
  colbacklower=black!3,
  segmentation style={solid,draw=black!25,line width=0.35pt},
  fontlower=\small\color{notegray},
  before lower={\noindent\textbf{Řešení:}\enspace},
  leftlower=3mm,rightlower=3mm
}}
\newtcolorbox{prikladReseniBox}[1][]{solutionlesson,title={Příklad},#1}
\newtcolorbox{zadaniReseniBox}[1][]{solutionlesson,title={Zadání},#1}

% Pomucky pro autora pripravy: sedive poznamky; \reseni pro starsi pripravy
\newcommand{\poznamka}[1]{\par\begingroup\small\color{notegray}\itshape
  \textnormal{\textbf{Poznámka pro učitele.}}\enspace #1\par\endgroup}
\newcommand{\reseni}[1]{\par\begingroup\small\color{notegray}
  \textbf{Očekávané řešení:}\enspace #1\par\endgroup}
\newcommand{\blok}[1]{\needspace{5\baselineskip}\subsection*{#1}}
\newcommand{\vyrok}[1]{\uv{#1}}

\usepackage{tikz}
\usetikzlibrary{arrows.meta}
\nastavPripravu{1. ročník}{Funkce}{61}{Monotónnost a funkce prostá, lichá a sudá}
\begin{document}
\poznamka{Základ 45 minut: připomenutí grafu a monotónnosti 4 min,
čtení intervalů 7 min, tři definice monotónnosti 7 min,
sudost a lichost 8 min, prostota 7 min, periodicita s příkladem 4 min,
samostatná práce 6 min, organizace 2 min.
Formální důkazy a náročnější úlohy jsou rezerva.
Souřadnice bodů zapisujeme $[x;y]$ a intervaly českým způsobem.}

\blok{Monotónnost: jak se funkce mění?}
U funkce už známe její definiční obor a obor hodnot.
Teď sledujeme, \emph{co se děje s funkční hodnotou, když argument roste}.
Graf čteme zleva doprava; zatím nemusí být přímkou.

\begin{center}
\begin{tikzpicture}[x=.78cm,y=.58cm,>=Latex]
\draw[step=1,gray!22,very thin] (-5,-3) grid (5,3);
\draw[->] (-5.15,0)--(5.3,0) node[right] {$x$};
\draw[->] (0,-3.12)--(0,3.27) node[above] {$y$};
\foreach \x in {-4,-3,-2,-1,1,2,3,4}{\draw (\x,.07)--(\x,-.07) node[below=2pt,font=\scriptsize] {$\x$};}
\foreach \y in {-2,-1,1,2}{\draw (.07,\y)--(-.07,\y) node[left=2pt,font=\scriptsize] {$\y$};}
\draw[very thick] (-4,0)--(-2,2)--(-1,2)--(1,-2)--(2,-2)--(4,0);
\fill (-4,0) circle (2.7pt);\draw[very thick,fill=white] (4,0) circle (2.8pt);
\end{tikzpicture}
\end{center}

Graf funkce $f$ má $D(f)=\langle-4;4)$. Například
$f(-3)=1$, $f(0)=0$ a $f(2)=-2$; hodnota $f(4)$ neexistuje.

\begin{definiceBox}
Funkce $f$ je na intervalu $I\subseteq D(f)$ \textbf{rostoucí}, jestliže
pro každá $x_1,x_2\in I$ platí
\[
x_1<x_2\ \Longrightarrow\ f(x_1)<f(x_2).
\]
\end{definiceBox}
\begin{definiceBox}
Funkce $f$ je na intervalu $I\subseteq D(f)$ \textbf{klesající}, jestliže
pro každá $x_1,x_2\in I$ platí
\[
x_1<x_2\ \Longrightarrow\ f(x_1)>f(x_2).
\]
\end{definiceBox}
\begin{definiceBox}
Funkce $f$ je na intervalu $I\subseteq D(f)$ \textbf{konstantní}, jestliže
pro všechna $x_1,x_2\in I$ platí $f(x_1)=f(x_2)$.
\end{definiceBox}

\begin{prikladReseniBox}[title={Příklad 1: čtení monotónnosti z grafu}]
Urči z úvodního grafu intervaly růstu, poklesu a konstantnosti.
Je funkce rostoucí na celém $D(f)$?
\tcblower
Roste na $\langle-4;-2\rangle$ a $\langle2;4)$;
klesá na $\langle-1;1\rangle$;
konstantní je na $\langle-2;-1\rangle$ a $\langle1;2\rangle$.
Na celém oboru neroste, protože například $-1<1$,
ale $f(-1)=2>f(1)=-2$.
\end{prikladReseniBox}

\poznamka{Rozlišuj „rostoucí“ od „kladná“: na úseku
$\langle2;4)$ funkce roste, ačkoliv její hodnoty jsou záporné.
Jde o přísnou monotónnost; konstantní úsek proto není rostoucí.
Pojmy neklesající a nerostoucí případně přidej jen jako rezervu.}

\blok{Sudá a lichá funkce: souměrnost grafu}
Nestačí porovnat jen dvě vybrané hodnoty. Nejdřív musí platit,
že s každým přípustným $x$ leží v definičním oboru i $-x$.

\begin{definiceBox}
Funkce $f$ je \textbf{sudá}, jestliže je její definiční obor
souměrný vzhledem k nule a pro každé $x\in D(f)$ platí
\[
f(-x)=f(x).
\]
Graf sudé funkce je souměrný podle osy $y$.
\end{definiceBox}
\begin{definiceBox}
Funkce $f$ je \textbf{lichá}, jestliže je její definiční obor
souměrný vzhledem k nule a pro každé $x\in D(f)$ platí
\[
f(-x)=-f(x).
\]
Graf liché funkce je středově souměrný podle počátku.
\end{definiceBox}

\begin{center}
\begin{minipage}{.43\linewidth}\centering
\begin{tikzpicture}[x=.64cm,y=.38cm,>=Latex]
\draw[->] (-2.4,0)--(2.4,0) node[right] {$x$};
\draw[->] (0,-.3)--(0,4.5) node[above] {$y$};
\draw[very thick,domain=-2:2,samples=60] plot (\x,{\x*\x});
\fill (-1,1) circle (2.5pt);\fill (1,1) circle (2.5pt);
\end{tikzpicture}
\par\small $f(x)=x^2$ — sudá
\end{minipage}\hfill
\begin{minipage}{.43\linewidth}\centering
\begin{tikzpicture}[x=.7cm,y=.48cm,>=Latex]
\draw[->] (-2.3,0)--(2.3,0) node[right] {$x$};
\draw[->] (0,-2.2)--(0,2.2) node[above] {$y$};
\draw[very thick,domain=-1.25:1.25,samples=60] plot (\x,{\x*\x*\x});
\fill (-1,-1) circle (2.5pt);\fill (1,1) circle (2.5pt);
\end{tikzpicture}
\par\small $g(x)=x^3$ — lichá
\end{minipage}
\end{center}

\begin{prikladReseniBox}[title={Příklad 2: rozhodnutí z předpisu}]
Rozhodni, zda jsou funkce $u(x)=x^2+2$ a $v(x)=x^3-x$
na $\mathbb R$ sudé, liché, nebo ani jedno.
\tcblower
$u(-x)=(-x)^2+2=x^2+2=u(x)$: funkce $u$ je sudá.
$v(-x)=(-x)^3-(-x)=-x^3+x=-v(x)$: funkce $v$ je lichá.
U obou je $D=\mathbb R$, tedy obor je souměrný vůči nule.
\end{prikladReseniBox}

\poznamka{Na omezeném oboru nestačí sudý tvar předpisu:
$x^2$ na $\langle0;2\rangle$ není sudá funkce,
protože její definiční obor není souměrný vzhledem k nule.}

\blok{Prostá funkce: různé vstupy, různé výstupy}
Sudost a lichost se ptají na vztah $f(-x)$ a $f(x)$.
\textbf{Prostota} je jiná vlastnost: zda mohou dva různé
argumenty dávat stejnou funkční hodnotu.

\begin{definiceBox}
Funkce $f$ je \textbf{prostá}, jestliže pro všechna
$x_1,x_2\in D(f)$ platí
\[
f(x_1)=f(x_2)\ \Longrightarrow\ x_1=x_2.
\]
Ekvivalentně: dvěma různým argumentům odpovídají dvě různé hodnoty.
\end{definiceBox}

\begin{prikladReseniBox}[title={Příklad 3: prostá, nebo neprostá?}]
Rozhodni o prostotě $a(x)=2x+1$ a $b(x)=x^2$ na $\mathbb R$.
Co se stane s $b$, omezíme-li její obor na $\langle0;+\infty)$?
\tcblower
$a$ je prostá: z $2x_1+1=2x_2+1$ plyne $x_1=x_2$.
$b$ na $\mathbb R$ není prostá, protože $b(-2)=b(2)=4$,
ačkoliv $-2\ne2$.
Na $\langle0;+\infty)$ už $b$ prostá je:
dvě různá nezáporná čísla mají různé druhé mocniny.
\end{prikladReseniBox}

\poznamka{Grafické pravidlo: prostou funkci každá vodorovná přímka
protne nejvýše jednou. Nezaměňovat se svislicovým testem,
kterým rozhodujeme, zda je křivka vůbec grafem funkce.}

\blok{Periodická funkce: opakování hodnot}
Periodická funkce opakuje své hodnoty po stále stejném posunutí argumentu.

\begin{definiceBox}
Funkce $f$ je \textbf{periodická}, jestliže existuje číslo $T>0$
takové, že pro každé $x\in D(f)$ platí $x-T,x+T\in D(f)$ a
\[
f(x+T)=f(x).
\]
Číslo $T$ nazýváme \textbf{periodou funkce}.
\end{definiceBox}

\begin{prikladReseniBox}[title={Příklad 4: periodická funkce}]
Ukaž, že funkce $f(x)=\sin x$ na $\mathbb R$ je periodická.
\tcblower
Pro každé $x\in\mathbb R$ platí $\sin(x+2\pi)=\sin x$.
Periodou je tedy $T=2\pi$; například
$f(0)=f(2\pi)=0$ a $f(\frac\pi2)=f(\frac{5\pi}2)=1$.
\end{prikladReseniBox}

\blok{Samostatná práce}
\begin{zadaniReseniBox}[title={Procvičování}]
\begin{enumerate}[label=\alph*)]
\item Rozhodni, zda $p(x)=-3x+1$ roste, nebo klesá na $\mathbb R$.
\item Rozhodni o sudosti či lichosti $q(x)=x^4-2x^2$
    a $r(x)=x^3+2x$ na $\mathbb R$.
\item Je funkce $s(x)=|x|$ prostá na $\mathbb R$?
\item Je $t(x)=x^2$ sudá na oboru $\langle-2;2\rangle$?
    A na oboru $\langle0;2\rangle$?
\end{enumerate}
\tcblower
\begin{enumerate}[label=\alph*)]
\item Klesá: $p(x_2)-p(x_1)=-3(x_2-x_1)<0$ pro $x_1<x_2$.
\item $q(-x)=q(x)$, takže $q$ je sudá;
    $r(-x)=-r(x)$, takže $r$ je lichá.
\item Není: $s(-1)=s(1)=1$, ale $-1\ne1$.
\item Na $\langle-2;2\rangle$ je sudá; na $\langle0;2\rangle$
    není sudá ani lichá, protože obor není souměrný vůči nule.
\end{enumerate}
\end{zadaniReseniBox}

\blok{Domácí procvičení}
\begin{zadaniReseniBox}[title={Domácí úkol}]
\begin{enumerate}[label=\alph*)]
\item Z úvodního grafu vypiš intervaly růstu, poklesu
    a konstantnosti.
\item Urči monotonii $f(x)=4-2x$ na $\mathbb R$.
\item Rozhodni o sudosti/lichosti $g(x)=x^5-3x$
    a $h(x)=x^2+x$ na $\mathbb R$.
\item Rozhodni o prostotě funkcí $p(x)=x^3$ a $q(x)=x^2$
    na $\mathbb R$.
\item Je $u(x)=x^2$ na $D(u)=\{-2,-1,1,2\}$ sudá
    a je prostá?
\end{enumerate}
\tcblower
\begin{enumerate}[label=\alph*)]
\item Roste na $\langle-4;-2\rangle$ a $\langle2;4)$,
    klesá na $\langle-1;1\rangle$,
    konstantní na $\langle-2;-1\rangle$ a $\langle1;2\rangle$.
\item Klesá, protože $f(x_2)-f(x_1)=-2(x_2-x_1)<0$ při $x_1<x_2$.
\item $g(-x)=-g(x)$, tedy lichá; $h(-x)=x^2-x$ není
    obecně $h(x)$ ani $-h(x)$, tedy ani sudá, ani lichá.
\item $p$ je prostá (třetí mocnina je přísně rostoucí),
    $q$ není prostá, protože $q(-1)=q(1)$.
\item Je sudá, protože obor je souměrný a $u(-x)=u(x)$;
    prostá není, např. $u(-2)=u(2)=4$.
\end{enumerate}
\end{zadaniReseniBox}
\end{document}
